\section{From sampling distributions to statistical inference}

The previous chapter introduced a second level of randomness. A probability model
describes individual observations, but it also determines the probability law of
the entire sample and of every statistic computed from that sample.

If
\[
X_1,X_2,\ldots,X_n
\]
is a random sample and
\[
T=g(X_1,X_2,\ldots,X_n)
\]
is a statistic, then \(T\) has a sampling distribution. Before the data are
observed, the statistic is random. After the sample has been collected, it takes
one concrete value
\[
t=g(x_1,x_2,\ldots,x_n).
\]

This chapter asks how the observed value \(t\) can be used to learn about the
unknown mechanism that generated the data.

\subsection*{Probability and inference move in opposite directions}

In a probability problem, the model is known. We begin with assumptions such as
\[
X_i\sim F
\]
or
\[
P(X_i=1)=p,
\]
and ask what kinds of samples and statistic values the model may produce.

The direction of reasoning is
\[
\text{model}
\longrightarrow
\text{possible samples}
\longrightarrow
\text{possible statistic values}.
\]

In statistical inference, the situation is reversed. We observe one sample and
one statistic value, while some part of the model is unknown. We then ask what
this observed result suggests about the mechanism that produced it.

The direction now appears to be
\[
\text{observed sample}
\longrightarrow
\text{information about the model}.
\]

This reversal is not automatic. One observed sample does not determine the model
uniquely. Different models may produce similar data, and even a correct model may
occasionally produce unusual samples. Statistical inference must therefore use
probability to measure uncertainty rather than pretend that uncertainty has
disappeared.

\begin{summarybox}
Probability asks what data a model may generate. Statistical inference asks what
an observed dataset allows us to say about the model that generated it.
\end{summarybox}

\subsection*{Parameters and statistics}

A probability model often contains an unknown numerical characteristic called a
\textbf{parameter}. We denote it by a symbol such as
\[
\theta,
\qquad
\mu,
\qquad
\sigma^2,
\qquad
\text{or}
\qquad
p.
\]
The parameter belongs to the model. It is not computed from the observed sample.

A statistic is computed from the sample and may be used to estimate the
parameter. We often write
\[
\widehat\theta=g(X_1,\ldots,X_n).
\]
Before observation, \(\widehat\theta\) is random. After observation, it becomes a
number
\[
\widehat\theta_{\mathrm{obs}}=g(x_1,\ldots,x_n).
\]

For example, in a Bernoulli model the unknown parameter is the success
probability \(p\). The sample proportion
\[
\widehat p=\frac1n\sum_{i=1}^nX_i
\]
is a statistic used to estimate it. Similarly, the sample mean
\[
\overline X=\frac1n\sum_{i=1}^nX_i
\]
may be used to estimate the population mean \(\mu\).

\begin{remarkbox}
A parameter and a statistic may play similar roles in a formula, but they are
mathematically different. The parameter is a fixed but unknown feature of the
model. The statistic is a random quantity before the sample is observed.
\end{remarkbox}

\subsection*{Why large samples matter}

The sampling distribution of a statistic may be complicated for small samples.
As the sample size increases, however, important regularities begin to appear.
Sample averages become more stable, their random fluctuations occur on a
predictable scale, and in many situations their standardized distributions take
an approximately universal form.

These regularities are the bridge from probability to inference. They allow us
to answer questions such as:

\begin{itemize}
    \item How close is a sample mean likely to be to the population mean?
    \item How much can a sample proportion vary from sample to sample?
    \item Which parameter values are compatible with the observed statistic?
    \item Would the observed result be plausible under a proposed model?
\end{itemize}

The first major regularity is the law of large numbers. It explains why averages
and proportions become stable. The second is the central limit theorem. It
describes the scale and approximate shape of their remaining fluctuations.

\subsection*{Inference is based on hypothetical repetition}

Usually we observe only one dataset. Nevertheless, statistical reasoning asks
what would happen if the same sampling procedure were repeated many times.

This does not mean that the actual experiment must literally be repeated. The
repetition is part of the probability model. It allows us to describe how the
statistic would vary across all samples that could have been observed.

Suppose a model with parameter \(\theta\) implies that a statistic \(T\) is
usually close to a certain value. After observing \(T=t_{\mathrm{obs}}\), we may
ask whether this result is typical or unusual under that model. If it is typical,
the data are compatible with the model. If it is extremely unusual, the data
provide evidence against it.

The same sampling behaviour can also be used in the opposite way. If a statistic
is usually close to the parameter it estimates, then an observed statistic can
be surrounded by a range of plausible parameter values.

These two ideas lead to the main forms of inference developed later in the
chapter:

\[
\text{estimation and confidence intervals},
\]

and

\[
\text{statistical hypotheses and tests}.
\]

\begin{theorembox}
Statistical inference is possible because statistics have predictable sampling
behaviour. Large-sample regularities connect the unknown model parameter with the
random statistic computed from observed data.
\end{theorembox}

\subsection*{The logical structure of the chapter}

The reasoning developed in this chapter follows the chain
\[
\text{probability model}
\longrightarrow
\text{sampling distribution}
\longrightarrow
\text{large-sample regularity}
\longrightarrow
\text{observed statistic}
\longrightarrow
\text{statistical conclusion}.
\]

The conclusion is never obtained from the observed number alone. It depends on
the sampling model, the assumptions, and the mathematical description of how the
statistic behaves across possible samples.

\begin{summarybox}
Statistical inference does not replace probability theory. It uses probability
theory in reverse. We first understand how a model generates samples and
statistics. Only then can one observed statistic be used to estimate an unknown
parameter or assess the compatibility of a proposed model with the data.
\end{summarybox}
